\documentclass{article}
\usepackage[a4paper, margin=2.5cm]{geometry}
\usepackage{parskip}

\begin{document}

\begin{center}
    {\large \textbf{ SURAT KUASA KHUSUS }} \\[0.5em]
    \rule{8cm}{0.4pt}
\end{center}

Yang bertanda tangan di bawah ini:
\begin{tabbing}
    \hspace{3cm} \= \hspace{1cm} \= \kill
    \textbf{Nama Lengkap} \> : \> Hendra Wijaya, S.E., M.Si. \\
    \textbf{Umur} \> : \> 57 Tahun \\
    \textbf{Pekerjaan} \> : \> Mantan Kepala Dinas PUPR Kota Tirta Jaya \\
    \textbf{Alamat} \> : \> Jl. Melati No. 45, Kota Tirta Jaya \\
    \textbf{NIK} \> : \> 3271041206690001
\end{tabbing}

Dalam hal ini bertindak untuk dan atas nama diri sendiri selaku \textbf{TERSANGKA} dalam perkara dugaan tindak pidana korupsi, selanjutnya disebut sebagai \textbf{PEMBERI KUASA};

\vspace{0.8em}
\begin{center}
    \textbf{DENGAN INI MEMBERIKAN KUASA HUKUM SEPENUHNYA KEPADA}
\end{center}
\vspace{0.3em}

\begin{tabbing}
    \hspace{3cm} \= \hspace{1cm} \= \kill
    \textbf{Nama Lengkap} \> : \> Liya, S.H. \\
    \textbf{Pekerjaan} \> : \> Advokat Penasihat Hukum \\
    \textbf{Alamat Kantor} \> : \> Jl. Jati No. 1, Kota Tirta Jaya \\
\end{tabbing}

Selanjutnya disebut sebagai \textbf{PENERIMA KUASA};

\vspace{0.8em}
\begin{center}
    \textbf{------------------------- KHUSUS -------------------------}
\end{center}
\vspace{0.3em}

Untuk dan atas nama Pemberi Kuasa, bertindak mendampingi, membela, serta memberikan bantuan hukum kepada Pemberi Kuasa yang sedang menghadapi proses hukum di tingkat Penyidikan, Penuntutan, hingga Persidangan di Pengadilan Negeri Tirta Jaya sehubungan dengan perkara dugaan Tindak Pidana Korupsi Proyek Pembangunan Sistem Pengolahan Air Bersih (SPAB) Kota Tirta Jaya Tahun Anggaran 2021 sebagaimana dimaksud dalam Pasal 603 dan/atau Pasal 604 Undang-Undang Nomor 1 Tahun 2023 tentang KUHP jo. Undang-Undang Nomor 1 Tahun 2026 tentang Penyesuaian Pidana jo. Undang-Undang Pemberantasan Tindak Pidana Korupsi.

Untuk keperluan tersebut, Penerima Kuasa dikuasakan untuk:
\begin{enumerate}
    \item Mendampingi Pemberi Kuasa pada setiap pemeriksaan di tingkat penyidikan, penuntutan, maupun persidangan di Pengadilan;
    \item Menghadap kepada Pejabat Penyidik, Penuntut Umum, Hakim, Majelis Hakim, serta instansi berwenang lainnya;
    \item Meminta turunan Berita Acara Pemeriksaan (BAP), mengajukan saksi/ahli yang meringankan, serta membela hak-hak hukum Pemberi Kuasa;
    \item Membuat, menandatangani, dan mengajukan eksepsi, pembelaan (pledoi), duplik, serta memorandum banding/kasasi apabila diperlukan;
    \item Melakukan segala tindakan hukum lain yang dianggap perlu demi kepentingan pembelaan hukum Pemberi Kuasa sesuai dengan Kitab Undang-Undang Hukum Acara Pidana (KUHAP) dan peraturan perundang-undangan yang berlaku.
\end{enumerate}

Demikian Surat Kuasa ini diberikan dengan hak substitusi, baik seluruhnya maupun sebagian, serta dengan tegas memberikan hak retensi menurut undang-undang. 
\vspace{1.2em}

\begin{flushright}
    Tirta Jaya, 17 September 2026 \\[2em]
\end{flushright}

\noindent
\begin{minipage}[t]{0.48\textwidth}
    Penerima Kuasa, \\[1em]
    % Kotak ilustrasi e-meterai darurat
    \framebox{
        \begin{minipage}{3.2cm}
            \centering
            \tiny \textbf{E-MATERAI} \\
            \vspace{0.2em}
            \scriptsize Rp10.000 \\
            \tiny [ Sepuluh ribu rupiah ]
        \end{minipage}
    }\\[0.5em]
    \textbf{Liya, S.H.}
\end{minipage}
\hfill
\begin{minipage}[t]{0.48\textwidth}
    Pemberi Kuasa, \\[3.5em]
    \textbf{HENDRA WIJAYA, S.E., M.Si.}
\end{minipage}

\end{document}
